% ============================================================
%  Chapitre 7 — Conclusion et Perspectives
% ============================================================
\chapter{Conclusion et Perspectives}
\thispagestyle{fancy}

\section{Synthèse des réalisations}

Ce projet de fin de cycle a abouti au développement complet de \farmAI{} v3.0,
une plateforme agro-technologique innovante répondant aux besoins spécifiques de
l'agriculture tunisienne. Les principaux livrables réalisés, sprint par sprint, sont :

\begin{table}[H]
\centering
\caption{Synthèse des livrables par sprint}
\label{tab:deliverables}
\begin{tabular}{clll}
\toprule
\textbf{Sprint} & \textbf{Titre} & \textbf{Livrables réalisés} & \textbf{Statut} \\
\midrule
S1 & Infra \& Auth &
  FastAPI 50+ endpoints, JWT+OTP, PostgreSQL 30+ tables,
  React/Vite scaffold & \yes~Livré \\
S2 & IoT \& Télémétrie &
  NODE\_A et NODE\_B firmware (HTTP POST), simulation Wokwi,
  dashboard télémétrie Recharts & \yes~Livré \\
S3 & Module Apiculture &
  HiveIntel 7 sous-modules, CRUD complet, intégration NODE\_B,
  QR Codes & \yes~Livré \\
S4 & Vision IA \& Scanner &
  8 modèles YOLO v11, Scanner PWA 13 modèles, Map SIG Leaflet,
  Expert FAB & \yes~Livré \\
S5 & Assistant Souverain &
  Labess-7B RAG Darija, STT/TTS, LLaVA vision, Lite Mode Oracle & \yes~Livré \\
S6 & MLOps \& Final &
  Pipeline DVC+MLflow, IsolationForest (Sil.=0.73),
  Docker Compose, Oracle Cloud & \yes~Livré \\
\bottomrule
\end{tabular}
\end{table}

\section{Contributions originales}

\farmAI{} v3.0 apporte trois contributions originales au domaine de l'agritech tunisienne :

\subsection{Souveraineté IA en Darija}

La principale contribution de ce projet est la mise en œuvre d'un assistant
conversationnel souverain en dialecte arabe tunisien (Darija). En utilisant le modèle
Labess-7B \cite{labess7b} couplé à un système RAG enrichi de sources officielles
tunisiennes (UTAP, AVFA), la plateforme offre aux agriculteurs un interlocuteur
numérique dans leur langue natale — sans dépendance à des API cloud étrangères.
Il s'agit d'une avancée significative vers la démocratisation de l'IA pour les
populations rurales tunisiennes.

\subsection{Intégration IoT-IA complète}

La chaîne complète ESP32 (Wokwi simulation → matériel réel) → FastAPI → IsolationForest
→ Dashboard → WebSocket constitue une architecture de référence pour les projets
IoT agricoles : même code firmware en simulation et en production, détection
d'anomalies automatique, alertes temps réel.

\subsection{MLOps agricole reproductible}

L'application rigoureuse des pratiques MLOps (DVC + MLflow) à un contexte agricole
tunisien démontre la faisabilité d'une industrialisation des modèles ML pour les
petites et moyennes exploitations, sans infrastructure GPU coûteuse.

\section{Difficultés rencontrées}

\begin{enumerate}
  \item \textbf{Déséquilibre des datasets YOLO :} Certaines classes de pathologies
        agricoles tunisiennes (maladies spécifiques à l'olivier, insectes ravageurs
        locaux) sont sous-représentées dans les datasets publics. Solution : stratégie
        d'augmentation aggressive (Mosaic, MixUp) et pondération des classes.

  \item \textbf{Latence LLM sur CPU :} Labess-7B sur CPU génère une réponse en 15-25
        secondes, impactant l'expérience utilisateur. Solution provisoire : Groq API
        comme fallback, et Lite Mode (base de connaissances statique) pour Oracle.

  \item \textbf{Wokwi et PlatformIO :} La simulation nécessite que le fichier ELF
        soit compilé avant de démarrer Wokwi. Sans le rechargement de VS Code après
        la première compilation, le bouton Play reste inactif.

  \item \textbf{Oracle Cloud ARM64 sans GPU :} Les modèles YOLO et Ollama ne fonctionnent
        pas sur l'instance ARM64 gratuite. Solution : LITE\_MODE=true désactive YOLO/LLM
        et active la base de connaissances statique.

  \item \textbf{Multi-tenancy et isolation :} Garantir l'isolation stricte des données
        entre fermes (même utilisateur ne peut pas voir les données d'un autre owner)
        a nécessité la propagation systématique du \code{farm\_id} dans toutes les
        requêtes et tous les modèles SQLAlchemy.
\end{enumerate}

\section{Perspectives d'évolution}

\begin{tcolorbox}[infobox, title={Feuille de route v4.0}]
\begin{description}
  \item[\textbf{Court terme (3-6 mois) :}]
    \begin{itemize}
      \item Intégration de microphones pour l'analyse audio des ruches
            (détection des fréquences 200-300 Hz caractéristiques de l'essaimage)
      \item Migration vers TimescaleDB pour les séries temporelles télémétriques
      \item Application React Native pour Worker (Android/iOS natif)
    \end{itemize}
  \item[\textbf{Moyen terme (6-18 mois) :}]
    \begin{itemize}
      \item ESP32-CAM pour photos automatiques des ruches toutes les heures
      \item Fine-tuning de Labess-7B sur un corpus agricole tunisien spécialisé
      \item Certification ISO 27001 pour la sécurité des données agricoles
      \item MLOps niveau 2 : CI/CD ML avec GitHub Actions
    \end{itemize}
  \item[\textbf{Long terme (18 mois+) :}]
    \begin{itemize}
      \item Réseau de fermes connectées pour l'analyse comparative régionale
      \item Partenariat UTAP/AVFA pour validation officielle des recommandations
      \item Dataset public tunisien de maladies agricoles (contribution open-source)
    \end{itemize}
\end{description}
\end{tcolorbox}

\section{Conclusion générale}

\farmAI{} v3.0 démontre qu'il est possible de construire une plateforme d'IA agricole
complète, souveraine et économiquement accessible en combinant des technologies
open-source de pointe (YOLO v11, FastAPI, ChromaDB, DVC, MLflow) avec une architecture
rigoureuse et des pratiques MLOps industrielles.

La spécificité tunisienne — assistant en Darija, données IoT adaptées aux conditions
climatiques locales, corpus RAG ancré dans les références officielles nationales — fait
de \farmAI{} plus qu'un prototype académique : c'est une fondation réelle pour la
transformation numérique de l'agriculture tunisienne.

\medskip
\begin{tcolorbox}[technote, title={Impact attendu}]
Si 1\% des 516 000 exploitants tunisiens adoptent \farmAI{} et réduisent leurs pertes
agricoles de 15\% grâce à la détection précoce des maladies, cela représente :
\textbf{774 fermes bénéficiaires} et une \textbf{économie estimée à 2.3 millions TND
par an} selon les données FAO sur les pertes moyennes par exploitation.
\end{tcolorbox}

Ce projet a été l'occasion d'appliquer l'ensemble des compétences acquises au cours
de la formation d'ingénieur en Data Science et Intelligence Artificielle à Tek-Up :
architecture logicielle, deep learning, MLOps, IoT, développement full-stack et
gestion de projet agile. Il constitue une contribution modeste mais sincère à la
transformation numérique de l'agriculture tunisienne.
